\chapter{Case Study and Data Engineering}
\label{chp:case_study_and_data}

This chapter documents the operational context in which the data-aware
simulation is discovered and the execution pipeline that turns raw
operational records into a simulation-ready event log. The chapter is
dedicated to only outline the design, therefore no models are fitted, no hypotheses are tested,
and no experimental parameters are chosen yet. Its prpose is to establish the
data foundation on which every subsequent claim in
Chapters~\ref{chp:experiments_and_framework}
and~\ref{chap:results} builds on, so that these later claims and results can be analysed against their exact input.

Section~\ref{sec:cs_environment} introduces the HHLA Container Terminal
Burchardkai (CTB) in Hamburg, Germany, the truck-handling scope of the study and the data
sources that feed the event-log engineering pipeline.
Section~\ref{sec:event_log_engineering} describes the multi-stage
transformation that produces the final event log, including
lifecycle-timestamp reconstruction and the derivation of case-level
context attributes required by the ProSiT simulation framework
\cite{Vinci2026prosit}. The chapter closes with a summary of the final
event-log characteristics that are used as input to the
train/test split defined in
Section~\ref{sec:train_test_split}.

%%%%%%%%%%%%%%%%%%%%%%%%%%%%%%%%%%%%%%%%%%%%%%%%%%%%%%%%%%%%%%
\section{Case Study Environment}
\label{sec:cs_environment}
%%%%%%%%%%%%%%%%%%%%%%%%%%%%%%%%%%%%%%%%%%%%%%%%%%%%%%%%%%%%%%

    \subsection{HHLA Container Terminal Burchardkai}

    The case studys envorinment is the HHLA Container Terminal Burchardkai
    (CTB), one of Europe's largest container terminals, located in the
    Port of Hamburg. CTB handles approximately 2.6 million TEU
    (Twenty-foot Equivalent Units) annually across 315 hectares
    \cite{hhla2025}.

    \textbf{Truck Process Scope:} External trucks visit CTB to deliver
    or receive containers. The process comprises:
    \begin{enumerate}
        \item \textbf{Gate In}: Truck registers at the entry gate.
        \item \textbf{Container handling}: Rail-Mounted Gantry cranes
              (RMGs), Van Carriers (VC) or manual operators (HO2) load
              or unload containers at storage blocks.
        \item \textbf{Gate Out}: Truck exits the terminal.
    \end{enumerate}

    The CTB environment in this project comes with several resource-constrains. Located between the waterside and the landside, 22 RMG storage blocks (T06--T27) are operating and serving both ends simulataniously with a maximum of 3 automatic Gantry Cranes in parallel. Additionally the CTB has two specialised areas. An empty-container depot (In operational terms also called "MT" or "LL" - in German for "Leerlager") and a manually operated yard area (In operational terms also called HO2). The Terminal Gate consists of 16 lanes half of which built for incoming trucks, the other half for outgoing trucks. The amount of operated gates per day and shift differs, depending on the resource planing by the terminal management. 
   
    \subsection{Data Sources}
    \label{sec:cs_data_sources}

    The operational data used in this thesis originates from several
    information systems that support day-to-day terminal operations.
    Historically, the information-system landscape at CTB evolved
    organically, resulting in a heterogeneous environment in which
    different systems manage distinct aspects of terminal operations,
    orchestrated by a Terminal Operating System. The primary data
    sources for this study are:

    \begin{itemize}
        \item \textbf{ITS schedule (scheduling system)} — records truck
              visit schedules, including timestamps for each operational
              step (gate-in time, ready time at a handling point, order
              fulfilment time, gate-out time).
        \item \textbf{Report system} — captures stage-level execution
              data, including gate entry times, exchange-area
              assignments and container transaction records, bridging
              yard operations and container master data.
        \item \textbf{Yard management system} — snapshots every four
              hours of all yard block capacities and utilisation rates.
        \item \textbf{Container master data} — container-level details
              (size, type, hazardous material, dwell time, inbound and
              outbound carrier).
    \end{itemize}

    These systems generate independent data pools with different
    identifier schemes, requiring the data-integration procedure
    described in Section~\ref{sec:event_log_engineering}. Historical
    data from March and April 2026 (8 weeks) was extracted, covering
    approximately 104,000 truck visits.

    \textbf{Data volume:}
    \begin{itemize}
        \item 113,466 container move events,
        \item 104,000 truck visits,
        \item 336 yard-state snapshots for (8 weeks $\times$ 7 days $\times$
              6 snapshots per day).
    \end{itemize}

%%%%%%%%%%%%%%%%%%%%%%%%%%%%%%%%%%%%%%%%%%%%%%%%%%%%%%%%%%%%%%%%%%%%%%%%%
\section{Event Log Engineering}
\label{sec:event_log_engineering}
%%%%%%%%%%%%%%%%%%%%%%%%%%%%%%%%%%%%%%%%%%%%%%%%%%%%%%%%%%%%%%%%%%%%%%%%%

Event-log engineering is the first stage of the experimental pipeline.
The objective is to construct an integrated, process-oriented event log
from the raw operational data that satisfies the requirements of both
process mining and data-aware simulation while retaining as much
operational context as possible. This section describes the pipeline
stage by stage.

\subsection{Raw Data Sources and Preparation}

\subsubsection{ITS Schedule Data}

The ITS scheduling system provides the core scheduling data for the
truck operations reconstructed in this thesis. Each record represents
one activity instance within a truck visit and carries the following
key attributes:

\begin{itemize}
    \item \texttt{FAHRPLAN\_UID} — unique identifier for the truck visit
          (case ID).
    \item \texttt{LKW\_KENNZEICHEN} — truck license plate number.
    \item \texttt{ANLAUFPUNKT\_HALTESTELLE} — exchange area identifier
          (e.g.\ T20, LL, VC).
    \item \texttt{FP\_ZEITPUNKT\_ERSTELLUNG} — timestamp when the
          schedule entry was created (gate-in time).
    \item \texttt{ALP\_ZEITPUNKT\_BEREITMELDUNG} — timestamp when the
          activity became ready for execution (enabled time).
    \item \texttt{ALP\_ZEITPUNKT\_ERFUELLUNG} — timestamp when the
          activity was completed.
    \item \texttt{FP\_ZEITPUNKT\_GATEOUT} — timestamp when the truck
          exited the terminal.
\end{itemize}

All timestamp fields are parsed into standardised \texttt{datetime}
objects using \texttt{pandas} with the format
\texttt{\%d.\%m.\%Y~\%H:\%M}. Invalid or missing timestamps are flagged
for removal during the cleaning phase.

\subsubsection{Report System Data}

The report system provides complementary stage-level execution data,
completing the details for gate-entry events, and it links each
container number to the specific truck visit at a yard handling point.
Key attributes are:

\begin{itemize}
    \item \texttt{TruckVisitGkey} — unique truck-visit identifier used
          for cross-system matching.
    \item \texttt{TruckLicenseNbr} — truck license plate.
    \item \texttt{CallupTime} — timestamp when the truck physically
          entered the terminal.
    \item \texttt{StageStartTime} — timestamp when the activity became
          ready for execution (enabled time).
    \item \texttt{StageId} — stage identifier
          (e.g.\ \texttt{OCR}, \texttt{Gate-In}, \texttt{yard},
          \texttt{Exit}).
    \item \texttt{TruckExchangeArea} — visited handling point in the
          yard for a specific container.
    \item \texttt{TranCtrNbr} — container number associated with the
          transaction.
    \item \texttt{UnitInboundCarrierId} — inbound carrier identifier
          (truck plate, vessel voyage number or train ID).
    \item \texttt{UnitOutboundCarrierId} — outbound carrier identifier
          (truck plate, vessel voyage number or train ID).
\end{itemize}

A container is considered ``picked up'' or ``delivered'' by a truck if
the corresponding inbound or outbound carrier identifier resolves to a
truck. The report system uses detailed handover-lane codes for the
exchange area; these must be mapped to the higher-level operational
zones used by the scheduling system. Specifically, all exchange areas
beginning with ``HP'' (\textit{Hochpalette}) are mapped to ``LL''
(\textit{Leerlager}) to align with the schedule system's terminology.
The mapping is implemented in
\texttt{002\_Prep\_area\_mapping\_report.py}.

\subsubsection{Container Master Data}

Container-level attributes are extracted from the container master data
and linked to events via the matched \texttt{container\_id}:

\begin{itemize}
    \item \textbf{Container size} — TEU equivalent (20', 40', 45').
    \item \textbf{Full/empty status} — binary indicator
          (\texttt{Leer}/\texttt{Voll}).
    \item \textbf{Reefer status} — binary indicator for refrigerated
          containers (derived from the ISO code).
    \item \textbf{Hazardous-material status} — binary indicator based on
          the IMCO code.
    \item \textbf{Container type} — standard, open-top, platform, bulk,
          etc.
    \item \textbf{Dwell time} — hours the container has already spent in
          the terminal (for receive operations).
    \item \textbf{Container complexity score} — composite metric
          combining size, full status, reefer, hazmat and dwell time to
          quantify handling difficulty.
\end{itemize}

The complexity score is derived during the container-master-data
preprocessing stage and combines container type, loading status,
hazardous-goods indicators, reefer requirements and other
handling-related properties into a single numerical measure. Higher
values indicate visits expected to require more complex handling.

\subsubsection{Yard State Data}

Yard state data provides time-series information on storage-block
capacities and utilisation rates. Each record contains
\texttt{ZEITPUNKT} (observation timestamp), \texttt{AREA} (storage area
identifier), \texttt{KAP\_TEU\_THEORETISCH} (theoretical TEU capacity),
\texttt{BEL\_TEU} (current TEU occupancy) and \texttt{AUSLASTUNG\_\%}
(utilisation percentage). Because observations are irregularly spaced,
the data is resampled to a uniform one-minute frequency using linear
interpolation on the numeric fields (script
\texttt{004\_Prep\_create\_interpolated\_yard\_state\_for\_s3.py}).

\subsubsection{Truck Demand Data}

Truck demand information is derived from the scheduling system by
aggregating truck arrivals over rolling time windows. For each exchange
area and one-minute time interval, the feature
\texttt{TRUCK\_DEMAND\_LAST\_1H} is calculated as the number of trucks
that became ready for processing in that area during the preceding
sixty minutes. The demand pipeline (script
\texttt{005\_Prep\_create\_truck\_demand.py}):

\begin{enumerate}
    \item Resamples the schedule data to one-minute frequency based on
          \texttt{ALP\_ZEITPUNKT\_BEREITMELDUNG}.
    \item Applies a sixty-minute rolling window to count truck arrivals
          per area.
    \item Produces a separate demand series for each operational area
          (Gate, T06--T27, VC, LL, MT).
    \item Aggregates the individual RMG blocks (T06--T27) into a single
          yard-section demand metric.
\end{enumerate}

The resulting demand features capture workload pressure at the time of
case arrival, which is hypothesised in
Section~\ref{sec:data_aware_hypotheses} to influence waiting times and
resource-allocation decisions.

\subsection{Data Cleaning and Integration}
\label{sec:data_cleaning}

\subsubsection{Cross-System Identifier Matching}

The scheduling and report systems use different identifiers for truck
visits. To create a unified event log they are linked through composite
keys. Truck visits are matched by
$(\texttt{TruckLicenseNbr}, \texttt{CallupTime})\equiv
(\texttt{LKW\_KENNZEICHEN}, \texttt{FP\_ZEITPUNKT\_ERSTELLUNG})$
using a left join that preserves all schedule records and augments them
with the report-system \texttt{TruckVisitGkey} when available. Container
numbers are then matched via
$(\texttt{TruckVisitGkey}, \texttt{TruckExchangeArea})\equiv
(\texttt{TruckVisitGkey}, \texttt{ANLAUFPUNKT\_HALTESTELLE})$.

For cases where the primary matching fails, a fallback strategy is
applied to single-exchange-area visits: if a truck visit contains only
one yard activity and no container identifier was matched, the
container ID is inferred from any other transaction associated with the
same \texttt{TruckVisitGkey}. The integration logic is implemented in
\texttt{01\_match\_container\_ids\_to\_schedule.py}.

\subsubsection{Invalid Process Instances}
\label{sec:invalid_process_instances}

The cleaning pipeline removes complete cases rather than individual
records: whenever a truck visit violates one of the defined quality
criteria the entire process instance is excluded, so that only
internally consistent process executions are retained. The primary
cleaning procedures are implemented in
\texttt{00\_its\_fahrplan\_raw\_to\_fahrplan\_cleaned\_0304.py}.

\paragraph{Missing lifecycle information.} Cases are removed if one or
more critical lifecycle timestamps are missing:
\texttt{FP\_ZEITPUNKT\_ERSTELLUNG},
\texttt{ALP\_ZEITPUNKT\_BEREITMELDUNG} or
\texttt{ALP\_ZEITPUNKT\_ERFUELLUNG}. Such records cannot be transformed
into valid lifecycle events. The quality report identified 1,267
affected cases ($\approx 1.34\%$ of all truck visits).

\paragraph{Temporal inconsistencies.} Cases are removed whenever
fundamental temporal constraints are violated. First, activities with
negative execution durations
($\texttt{ALP\_ZEITPUNKT\_ERFUELLUNG}<\texttt{ALP\_ZEITPUNKT\_BEREITMELDUNG}$)
indicate timestamp errors and are excluded. Second, overlapping
activities within a single visit ($\text{start}_{i+1}<\text{complete}_{i}$
after sorting) contradict the operational logic (a truck can only be at
one handling location at a time). The cleaning procedure identified
three cases with negative durations and 43 with overlapping activities.

\paragraph{Extreme runtime outliers.} Two further runtime filters are
applied. The first removes cases with excessive gate-out delay
($t_{\text{GateOut}}-t_{\text{LastCompletion}}>10\text{h}$); such
records commonly indicate system failure. The second removes visits
whose total turnaround time exceeds twelve hours
($t_{\text{GateOut}}-t_{\text{GateIn}}>12\text{h}$), which is
operationally implausible. The filters identified 22 extreme gate-out
cases and 22 twelve-hour outliers.

\subsubsection{Process Completeness Filters}

\paragraph{Container identification coverage.} Overall, 111,646 of
113,136 operational records could be assigned a valid
\texttt{TruckVisitGkey} (98.68\%), and 110,866 records could be assigned
a valid container identifier (97.99\%). Despite the fallback procedure,
2,270 records belonging to 2,034 truck visits remained without a valid
container assignment. Since container-related attributes are central to
the later context enrichment and simulation-discovery stages, these
process instances are removed entirely (a complete-case filter). This
removes 2,034 truck visits and 3,397 associated records, leaving 91,335
complete truck visits and 109,739 operational records.

\paragraph{Unknown activity types.} During activity construction, cases
containing activities classified as \texttt{unknown} are removed. Such
events cannot be reliably mapped to operational process steps and would
introduce ambiguity into both process discovery and simulation
discovery.

\paragraph{Process scope restriction.} The event log is restricted to
truck visits containing exactly one interaction with an automated RMG
block. Visits may still contain additional activities in other yard
areas (empty-container yard or VC yard), but only a single RMG block
interaction is permitted. This restriction ensures that RMG-specific
workload, utilisation and demand measures can be assigned
unambiguously at case level. It also forms the basis for the
block-focused what-if experiments in
Section~\ref{sec:what_if_design}.

\paragraph{Event-level cleanup.} Immediately before exporting the event
log, remaining events with an \texttt{\_unknown} suffix are removed.
Unlike the previous filters, this is an event-level cleanup rather than
a case-level exclusion; its purpose is to eliminate residual events
that cannot be interpreted reliably within the final process model.

\subsection{Case Definition}
\label{sec:case_definition}

A \textit{case} in the event log corresponds to a single truck visit,
uniquely identified by \texttt{FAHRPLAN\_UID}. Each case represents the
complete sequence of activities executed by one truck between gate
entry and gate exit.

The data-aware simulation framework used in this thesis
\cite{Vinci2026prosit} accepts only case-level contextual attributes.
To attach block-specific context features
(\texttt{target\_utilization}, \texttt{target\_demand},
\texttt{target\_rank}) consistently, the design decision was made to
include only process instances that interact with exactly one RMG block.
Visits involving multiple RMG blocks are excluded because block-specific
congestion measures could not be assigned unambiguously. Visits to
other yard areas (VC or LL) that also involve exactly one RMG block are
retained. The filtering logic is implemented in
\texttt{filter\_one\_rmg\_only.py}.

\subsection{Activity Construction}
\label{sec:activity_construction}

Activities are defined by combining an \textit{operation type}
(receive, delivery or mixed) with a \textit{resource type} (an RMG
block, LL or HO2). Table~\ref{tab:activities} lists the resulting
labels.

\begin{table}[h]
\centering
\caption{Activity labels used in the final event log.}
\label{tab:activities}
\begin{tabular}{ll}
\toprule
\textbf{Activity} & \textbf{Definition} \\
\midrule
\texttt{Gate In}       & Truck registration at entry gate \\
\texttt{RMG\_receive}  & RMG unloads container from truck (import) \\
\texttt{RMG\_delivery} & RMG loads container onto truck (export) \\
\texttt{RMG\_mixed}    & Both receive and delivery at same block \\
\texttt{LL\_receive}   & Empty-container depot: unload \\
\texttt{LL\_delivery}  & Empty-container depot: load \\
\texttt{LL\_mixed}     & Empty-container depot: unload and load \\
\texttt{HO2\_receive}  & Manual area: unload \\
\texttt{HO2\_delivery} & Manual area: load \\
\texttt{HO2\_mixed}    & Manual area: unload and load \\
\texttt{Gate Out}      & Truck departure \\
\bottomrule
\end{tabular}
\end{table}

The 22 individual RMG blocks (T06--T27) are retained as distinct
resources to preserve block-specific utilisation patterns.

\subsection{Activity Timestamps and Temporal Semantics}
\label{sec:lifecycle_events}

The CTB operational systems provide timestamps from which two observed
lifecycle states are constructed for each activity:

\begin{itemize}
    \item \textbf{Start time} (\texttt{start:timestamp}) — the moment
          execution began, mapped from the scheduling or report system.
    \item \textbf{Complete time} (\texttt{time:timestamp}) — the moment
          execution finished.
\end{itemize}

The \textit{service time} (or execution time) is the directly
identifiable temporal quantity:
\begin{equation}
    ET_i = t_{\mathrm{complete},i} - t_{\mathrm{start},i}.
\end{equation}

An explicit \texttt{enabled:timestamp} is \textbf{not} reconstructed
in the input event log.  Consistent with ProSiT
\cite{Vinci2026prosit} and earlier simulation-discovery work
\cite{Camargo2020,PMRozinat2009}, activity enablement is obtained from
process-model alignment during the simulation-parameter discovery
phase (Section~\ref{sec:simulation_discovery_methodology}).  ProSiT
uses A* alignment to determine when each transition becomes enabled in
the discovered Petri net, then computes a model-derived pre-service
delay relative to both the enablement time and the resource's
availability schedule.

\subsubsection{Timestamp Mapping by Activity Type}

\begin{itemize}
    \item \textbf{Gate In.}
          $\text{start}=\text{FP\_ZEITPUNKT\_ERSTELLUNG}$,
          $\text{complete}\approx\text{start}$.
          The OCR gate timestamp (\texttt{StageStartTime}) from the
          report system is retained as the auxiliary column
          \texttt{ocr\_timestamp} for diagnostic purposes but is not
          mapped to an enablement concept.
    \item \textbf{Yard activities.}
          $\text{start}=\text{ALP\_ZEITPUNKT\_BEREITMELDUNG}$,
          $\text{complete}=\text{ALP\_ZEITPUNKT\_ERFUELLUNG}$.
    \item \textbf{Gate Out.}
          $\text{start}=\text{complete}=\text{FP\_ZEITPUNKT\_GATEOUT}$.
\end{itemize}

\subsubsection{Model-Derived Pre-Service Delay}

Because activity enablement is derived from the process model, the
resulting quantity
\begin{equation}
    D_{\mathrm{pre\text{-}service},i}
    = t_{\mathrm{start},i} - t_{\mathrm{enabled,model},i}
\end{equation}
is a \textit{model-derived pre-service delay}.  In ProSiT terminology
this is called ``waiting time,'' but in the CTB application it must
not be interpreted as pure queueing or resource-induced waiting.
Physical truck movement between consecutive terminal locations is not
represented as a separate activity in the discovered process model.
Therefore the model-derived interval may include:
\begin{itemize}
    \item physical movement between operational locations,
    \item internal traffic and congestion effects,
    \item resource availability and queue waiting,
    \item other delays not explicitly represented in the control-flow
          model.
\end{itemize}
This is a domain-specific limitation when transferring
business-process simulation semantics to physical terminal operations
and constitutes an important methodological boundary of this thesis
(discussed further in Section~\ref{sec:limitations}).

\subsubsection{Methodological Note: The Abandoned Transition-Baseline Approach}
\label{sec:abandoned_baseline}

An earlier version of the event-log pipeline estimated lower-tail
route-specific transition durations from the scheduling data (5th
percentile of the empirical distribution per transition pair,
\texttt{transition\_baseline.csv}) and used these to reconstruct an
explicit \texttt{enabled:timestamp} in the input event log.  Closer
inspection revealed two issues:
\begin{enumerate}
    \item The statistical route baselines are not event-specific
          observations of physical arrival; they are aggregate quantiles
          that vary substantially with the chosen percentile
          (Section~\ref{sec:results_percentile_sens}).
    \item In some cases the estimated readiness timestamp occurs after
          the observed start timestamp, producing negative waiting
          times that had to be clipped to zero.
\end{enumerate}
Furthermore, ProSiT does not require an externally constructed
enabled timestamp — it derives enablement internally from the aligned
process model \cite{Vinci2026prosit}.  The final methodology therefore
retains only the observed start and completion timestamps and delegates
enablement reconstruction to the process-model alignment, improving
consistency between data semantics and simulation semantics.

The transition-baseline data is preserved in the repository as
diagnostic material and may be useful for future work that introduces
an explicit spatial-transition model.

\subsection{Context Enrichment}
\label{sec:context_enrichment}

To enable data-aware simulation, the event log is enriched with
contextual attributes describing the operational state at the time of
process execution. These attributes are organised into five
categories.

\subsubsection{Container Features}

Container-related context is aggregated to case level so that visit
properties can serve as case-level attributes in ProSiT (script
\texttt{03\_aggregate\_container\_details\_case\_level.py}):

\begin{itemize}
    \item \textbf{\texttt{n\_containers}} — number of unique containers
          handled during the visit.
    \item \textbf{\texttt{has\_hazardous}} — at least one IMDG container
          involved.
    \item \textbf{\texttt{has\_reefer}} — at least one reefer container
          involved.
    \item \textbf{\texttt{full\_ratio}} — share of loaded containers.
    \item \textbf{\texttt{container\_complexity\_mean}} — mean container
          complexity score.
    \item \textbf{\texttt{container\_complexity\_sum}} — cumulative
          complexity score.
\end{itemize}

Binary indicators are aggregated using logical OR, numerical measures
through summary statistics.

\subsubsection{Yard Features}

Yard state features capture the storage capacity and utilisation status
of different terminal areas at truck arrival (Gate In event). These
features are matched temporally using \texttt{pd.merge\_asof} with a
\texttt{nearest} strategy:

\begin{itemize}
    \item \texttt{gate\_utilization} — overall terminal-yard utilisation.
    \item \texttt{rmg\_utilization} — aggregate utilisation across all
          RMG blocks.
    \item \texttt{vc\_utilization} — VC-area utilisation.
    \item \texttt{mt\_utilization} — empty-container storage
          utilisation.
    \item \texttt{target\_utilization} — utilisation of the specific
          block assigned to the truck.
\end{itemize}

These features are added in \texttt{06\_add\_utils\_to\_eventlog.py}.

\subsubsection{Demand Features}

Demand features quantify operational load pressure at different
terminal areas over a rolling 60-minute window (scripts
\texttt{05\_add\_demand\_features\_to\_eventlog.py} and
\texttt{08\_add\_all\_target\_features.py}):

\begin{itemize}
    \item \texttt{gate\_demand} — trucks processed in the last hour
          across the terminal.
    \item \texttt{rmg\_demand} — trucks assigned to RMG blocks in the
          last hour.
    \item \texttt{vc\_demand} — demand for VC operations.
    \item \texttt{mt\_demand} — demand for empty-container handling.
    \item \texttt{target\_demand} — demand for the assigned RMG block.
\end{itemize}

\subsubsection{Target Block Features}

For single-block truck visits the assigned target RMG block is a
critical attribute. Beyond the target-specific utilisation and demand,
the log carries:

\begin{itemize}
    \item \texttt{target\_rank} — ordinal ranking of RMG blocks by
          real-time utilisation at case arrival.
    \item \texttt{target\_rank\_group} — categorised rank
          (\texttt{top\_3}, \texttt{mid}, \texttt{low}).
    \item \texttt{primary\_target\_area} / \texttt{target\_area} — the
          block identifier assigned to the case.
\end{itemize}

These features are extracted from the \texttt{org:resource} attribute
of RMG activities (\texttt{08\_add\_all\_target\_features.py}).

\subsubsection{Temporal Features Required by ProSiT}
\label{sec:temporal_features}

ProSiT 1.0.3 does not derive temporal features (hour of day, day of
week) directly from event timestamps during simulation-parameter
discovery. Time-dependent case attributes must be provided explicitly at
case level. Two features are therefore added at the case level:

\begin{itemize}
    \item \texttt{@@hour} — hour of day (0--23) extracted from the
          earliest \texttt{start:timestamp} of the case.
    \item \texttt{@@weekday} — day of week (0=Mon, 6=Sun).
\end{itemize}

The \texttt{@@} prefix indicates that these are case-level attributes
used by ProSiT. The temporal-feature extraction is performed in
\texttt{09\_eventlog\_to\_xes.py}. This explicit encoding of the
temporal signal is not an option in ProSiT but a requirement, and it
is one of the observed side effects of the ProSiT feature-discovery
implementation that are consolidated in
Section~\ref{sec:prosit_workarounds}.

\subsection{Final Event Log Structure}
\label{sec:eventlog_structure}

The final event log combines process-execution data, lifecycle
information and operational context within a single simulation-ready
representation. Table~\ref{tab:core_attrs} lists the core process
attributes;
Tables~\ref{tab:case_attrs}--\ref{tab:temporal_attrs} list the case-level
context attributes attached to every event.

\begin{table}[H]
\centering
\caption{Core process attributes.}
\label{tab:core_attrs}
\begin{tabular}{p{4cm} p{9cm}}
\hline
\textbf{Attribute} & \textbf{Description} \\
\hline
\texttt{case:concept:name} & Unique truck-visit identifier
                             (\texttt{FAHRPLAN\_UID}). \\
\texttt{concept:name}      & Activity executed during the visit. \\
\texttt{start:timestamp}   & Observed/reconstructed activity start timestamp. \\
\texttt{time:timestamp}    & Observed/reconstructed activity completion timestamp. \\
\texttt{org:resource}      & Operational location of the activity. \\
\hline
\end{tabular}
\end{table}

\begin{table}[h]
\centering
\caption{Case-level container and process-shape attributes.}
\label{tab:case_attrs}
\begin{tabular}{p{4cm} p{9cm}}
\hline
\textbf{Attribute} & \textbf{Description} \\
\hline
\texttt{process\_flow\_type} & Visit type (delivery, receive, mixed). \\
\texttt{n\_containers}       & Number of containers in the visit. \\
\texttt{n\_stops}            & Number of operational stops. \\
\texttt{n\_deliveries}       & Number of delivery transactions. \\
\texttt{n\_receives}         & Number of receive transactions. \\
\texttt{has\_hazardous}      & IMDG indicator. \\
\texttt{has\_reefer}         & Reefer indicator. \\
\texttt{full\_ratio}         & Share of loaded containers. \\
\texttt{visit\_complexity}   & Composite complexity score. \\
\hline
\end{tabular}
\end{table}

\begin{table}[h]
\centering
\caption{Demand and utilisation features at case arrival.}
\label{tab:util_attrs}
\begin{tabular}{p{4cm} p{9cm}}
\hline
\textbf{Attribute} & \textbf{Description} \\
\hline
\texttt{gate\_demand}       & Current gate demand. \\
\texttt{rmg\_demand}        & Current demand in the RMG yard. \\
\texttt{vc\_demand}         & Current demand in the VC area. \\
\texttt{mt\_demand}         & Current demand in the manual yard. \\
\texttt{gate\_utilization}  & Current gate utilisation. \\
\texttt{rmg\_utilization}   & Current RMG-yard utilisation. \\
\texttt{vc\_utilization}    & Current VC-area utilisation. \\
\texttt{mt\_utilization}    & Current manual-yard utilisation. \\
\hline
\end{tabular}
\end{table}

\begin{table}[h]
\centering
\caption{Target-area context attributes.}
\label{tab:target_attrs}
\begin{tabular}{p{4cm} p{9cm}}
\hline
\textbf{Attribute} & \textbf{Description} \\
\hline
\texttt{target\_utilization} & Utilisation of the primary target area
                                at truck arrival. \\
\texttt{target\_demand}      & Demand for the primary target area. \\
\texttt{target\_rank}        & Utilisation rank among all RMG blocks. \\
\texttt{target\_rank\_group} & Categorised rank
                                (\texttt{top\_3}, \texttt{mid},
                                \texttt{low}). \\
\hline
\end{tabular}
\end{table}

\begin{table}[h]
\centering
\caption{Temporal context attributes required by ProSiT.}
\label{tab:temporal_attrs}
\begin{tabular}{p{4cm} p{9cm}}
\hline
\textbf{Attribute} & \textbf{Description} \\
\hline
\texttt{@@hour}    & Hour of day derived from the case-start timestamp. \\
\texttt{@@weekday} & Day of week derived from the case-start timestamp. \\
\hline
\end{tabular}
\end{table}

All contextual attributes are attached at case level and therefore
remain constant throughout the case. This design is consistent with the
current ProSiT feature-discovery approach, which primarily utilises
case-level context information rather than event-level dynamic state
\cite{Vinci2026prosit}.

\subsection{Event Log Characteristics}
\label{sec:eventlog_characteristics}

After the full preprocessing pipeline the final event log contains
89,458 truck visits represented by 277,423 events. The activity
distribution is shown in Table~\ref{tab:activity_distribution}.
Table~\ref{tab:eventlog_stats} summarises the overall event log.

\begin{table}[h]
\centering
\caption{Activity distribution in the final event log.}
\label{tab:activity_distribution}
\begin{tabular}{lr}
\toprule
\textbf{Activity} & \textbf{Frequency} \\
\midrule
Gate In       & 89,458 \\
Gate Out      & 89,458 \\
RMG\_receive  & 38,334 \\
RMG\_delivery & 17,524 \\
RMG\_mixed    & 13,522 \\
LL\_delivery  & 9,589  \\
HO2\_delivery & 6,689  \\
HO2\_receive  & 4,762  \\
LL\_mixed     & 4,451  \\
HO2\_mixed    & 2,254  \\
LL\_receive   & 910    \\
\bottomrule
\end{tabular}
\end{table}

\begin{table}[h]
\centering
\caption{Final event log summary.}
\label{tab:eventlog_stats}
\begin{tabular}{lr}
\toprule
\textbf{Statistic} & \textbf{Value} \\
\midrule
Cases                  & 89,458 \\
Events                 & 277,423 \\
Distinct activities    & 11 \\
Observation start      & 2026-03-02 05:25 \\
Observation end        & 2026-05-01 (approx.) \\
Distinct RMG resources & 22 (T06--T27) \\
Additional pools       & LL, HO2, VC, Gate In, Gate Out \\
\bottomrule
\end{tabular}
\end{table}

Every case contains a \texttt{Gate In} event, one or more operational
yard activities and a final \texttt{Gate Out} event. Operational
activities may occur at an RMG block, the LL area, HO2 or the VC area,
depending on the operational requirements of the visit.

The generated event log is exported both as CSV
(\texttt{s6\_eventlog\_target\_rank\_features.csv}) and XES
(\texttt{s6\_eventlog\_target\_rank\_features.xes}). The CSV
representation serves as the primary format for the pipeline; the XES
representation enables compatibility with pm4py, ProM and ProSiT and
serves as the foundation for the process discovery in
Chapter~\ref{chp:experiments_and_framework}. The final XES export
contains 89,458 cases and 277,423 events.

The chronological character of the log is preserved end to end: every
case carries the same monotone timeline from Gate In to Gate Out. This
property is exploited in Chapter~\ref{chp:experiments_and_framework} to
define a temporal train/test split at case level (see
Section~\ref{sec:train_test_split}).
